\section{成像位置、清晰观察与视角放大率}

\subsection{逐段成像与距离符号}

本讲把折射、反射和薄透镜组合起来。每到一个新元件，都先在图上判断入射光正在发散还是会聚，再确定本次物距。沿该段入射光的传播方向，实物距 \(u>0\)，虚物距 \(u<0\)；沿出射方向，实像距 \(v>0\)，虚像距 \(v<0\)。反射会改变传播方向，下一段须重新量距。

同一外界介质中的薄透镜满足 \(1/u+1/v=1/f\)，横向放大率为 \(\beta=-v/u\)。会聚镜取 \(f>0\)，发散镜取 \(f<0\)。球面镜的近轴公式形式相同，凹面镜 \(f=R/2>0\)，凸面镜 \(f=-R/2<0\)。折射率不同的两侧应使用
\[
 \frac{n_1}{u}+\frac{n_2}{v}=\Phi,\qquad
 \beta=-\frac{n_1v}{n_2u}.
\]
薄平行玻璃板没有光焦度，却会改变后续光组的等效物距。法向近轴传播时，厚度 \(t\)、折射率 \(n\) 的玻璃段，在两侧空气的等效成像距离中贡献 \(t/n\)，不是 \(nt\)。\(nt\) 是法向光程，两种量不能混用。

\subsection{平行板的约化传播与光程}

图\ref{o4:plate-propagation}中，光从空气进入玻璃再回到空气，出射方向与入射方向平行。若空气中小倾角为 $\theta$，玻璃中为 $\theta/n$，穿过厚度 $t$ 时的横向增量为 $t\theta/n$。因此在成像的高度—方向传递中，这段相当于空气传播 $t/n$；若计算法向传播时间，光程却是 $nt$，用时为 $nt/c$。
\fig{o4:plate-propagation}{平行板后的出射线向后延长到较靠近观察侧的 $S_1$；约化传播长度与光程用于不同问题}{\begin{tikzpicture}

\fill[blue!10](0,-.3)rectangle(2,1.5);\draw[world](0,-.3)--(0,1.5) (2,-.3)--(2,1.5);\draw[axisline](-2.5,0)--(4.5,0);\draw[ray,blue](-2,0)--(0,.5);\draw[ray,blue](0,.5)--(2,.83333);\draw[ray,red](2,.83333)--(4.2,1.38333);\draw[virtual,red](-1.33333,0)--(2,.83333);\fill(-2,0)circle(1.5pt)node[below]{$S$};\fill(-1.33333,0)circle(1.5pt)node[below]{$S_1$};\node at(1,1.8){玻璃 $n$};\node at(-1.3,1.1){$\theta$};\node at(1,.35){$\theta/n$};\draw[dim](0,-.7)--(2,-.7)node[midway,below]{$t$};\draw[dim](-2,-1.3)--(-1.33333,-1.3)node[midway,below]{$\delta$};

\end{tikzpicture}}
将出射直线延长回轴，物点的表观移动为 $\delta=t-t/n$，方向朝出射侧。接入后续透镜时，应从透镜量到 $S_1$，而不是把光程 $nt$ 当作物距的一部分。板的两面平行、外侧介质相同和近法向传播共同保证了这个简化；棱镜折转还须先把反射后的路径展开，不能直接把所有空间距离按纸面方向相加。

\subsection{往返系统的反射与展开}

同一透镜在往返途中通常作用两次。图\ref{o4:roundtrip-unfold}左图逐次标出透射、反射、再次透射；右图把镜后区域展开，使反射后的左行段成为继续右行的段，回程透镜成为一片镜像透镜。展开保持长度和横向高度，改变的是轴向量距方式。
\fig{o4:roundtrip-unfold}{往返光路与展开光路：平面镜到透镜的距离在展开后累计两次，回程像也要映到展开侧}{\begin{tikzpicture}

\begin{scope}[scale=.85]
\draw[axisline](-1.7,0)--(2.5,0);\draw[world](0,-.85)--(0,.85) (2,-.85)--(2,.85);\fill[pattern=north east lines](2,-.85)rectangle(2.15,.85);\draw[axis,blue](-1,0)--(-1,.2)node[above]{$B$};\draw[axis,red](-1,0)--(-1,-.2)node[below]{$B'$};\draw[ray](-1,.2)--(0,.2);\draw[ray](0,.2)--(2,-.2);\draw[ray](2,-.2)--(0,-.6);\draw[ray](0,-.6)--(-1,-.2);\node[above]at(0,.85){透镜};\node[above]at(2,.85){平面镜};\draw[dim](0,-1.3)--(2,-1.3)node[midway,below]{$L$};\end{scope}
\begin{scope}[xshift=5cm,scale=.85]
\draw[axisline](-1.7,0)--(5.6,0);\draw[world](0,-.85)--(0,.85) (4,-.85)--(4,.85);\draw[normal](2,-.85)--(2,.85);\draw[axis,blue](-1,0)--(-1,.2)node[above]{$B$};\draw[axis,red](5,0)--(5,-.2)node[below]{$B'_{\rm unfolded}$};\draw[ray](-1,.2)--(0,.2);\draw[ray](0,.2)--(4,-.6);\draw[ray](4,-.6)--(5,-.2);\node[above]at(0,.85){去程透镜};\node[above]at(4,.85){回程镜像};\draw[dim](0,-1.3)--(4,-1.3)node[midway,below]{$2L$};\end{scope}

\end{tikzpicture}}
因此矩阵追迹时，平面镜与透镜间的两段传播可以合成 $2L$，但两次透镜作用都要保留。若银层在透镜某一曲面上，应先辨认实际经过的界面：平面镀银与凸面镀银产生不同的折射、反射次序。不能在未追迹前，把所有镀银透镜统一看成同一种等效镜。


\subsection{视角放大率与观察条件}

高度为 $h$的小物体放在眼前比较距离 $D$处时，裸眼小视角为 $\alpha_0\simeq h/D$。常取 $D=25\,\mathrm{cm}$作为放大率的约定基准；它不等于每个人的最近清晰距离。经光具观察的有向视角为 $\alpha$时，定义 $M=\alpha/\alpha_0$。负号表示倒像，$|M|$才是放大倍数；横向放大率 $h'/h$比较长度，与眼处视角一般不同。

以焦距 $f>0$的放大镜为例，图\ref{fig:o4-eye-distance}设眼睛 $E$ 在透镜 $O$ 后 $\ell$ 处，虚像到镜的距离大小为 $q$，像高为 $h'$。眼到像的距离是 $q+\ell$，所以小视角下
\[
 \alpha\simeq\frac{h'}{q+\ell},\qquad
 \alpha_0\simeq\frac hD,\qquad
 M=\frac{D}{q+\ell}\frac{h'}h.
\]
\fig{fig:o4-eye-distance}{眼到虚像的距离从眼睛量取，透镜物像距从光心量取；横纵尺度不同}{
\begin{tikzpicture}[x=.9cm,y=1cm]
\draw[axisline] (-5.6,0)--(2,0);
\draw[<->,>=Stealth,world] (0,-.5)--(0,.8);
\draw[axis,thick] (-.833333,0)--(-.833333,.25) node[above left]{$B$};
\draw[axis,thick] (-5,0)--(-5,1.5) node[above]{$B'$};
\draw[ray] (-.833333,.25)--(0,.25);\draw[ray] (0,.25)--(1.4,-.1);
\draw[ray] (-.833333,.25)--(0,0);\draw[ray] (0,0)--(1.4,-.42);
\draw[virtual] (-5,1.5)--(0,.25);\draw[virtual] (-5,1.5)--(0,0);
\draw[very thick] (1,-.4)--(1,.2);\node[above] at (1,.3){眼 $E$};\node[below left] at (0,0){$O$};
\draw[dim] (-5,-.9)--(0,-.9) node[midway,below]{$q$};
\draw[dim] (0,-.9)--(1,-.9) node[midway,below]{$\ell$};
\draw[dim] (-5,-1.6)--(1,-1.6) node[midway,below]{$q+\ell$};
\end{tikzpicture}}
比较距离和实际近点应先分清。设眼的实际近点为 $D_N$，它要求 $q+\ell=D_N$，而裸眼参照角仍为 $h/D$。联立透镜式可得更一般的近点视角放大率
\[
M_N=\frac{D}{D_N}\left(1+\frac{D_N-\ell}{f}\right).
\]
以下采用实际近点等于所选比较距离的通常模型，即 $D_N=D$。


近点观察要求眼到虚像恰为明视距离，故 $q+\ell=D$。再由 $v=-q$、$1/u-1/q=1/f$，得
\[
 u=\frac{fq}{q+f},\qquad \frac{h'}h=\frac qu,
 \qquad \boxed{M_D=1+\frac{D-\ell}{f}}\quad(0\le\ell<D).
\]
眼睛紧贴透镜时，才退回 $M_D=1+D/f$。无穷远观察时物在前焦面，出射倾角为 $h/f$，在光束确能进入瞳孔的条件下有 $M_\infty=D/f$。这里的 $D$ 始终是裸眼参照距离，不能因末像在无穷远就把分母参照角误取为零。有限孔径引起的遮光影响可见范围，不应在未经分析时直接改写成像公式。

放大率也不等于分辨能力。把已经模糊的像放得更大，不能恢复被衍射或像差丢失的细节；瞳孔与出瞳不匹配时，还可能只看见视场的一部分。因此分析仪器时应分别说明像在哪里、眼能否清晰调节、视角怎样改变，以及允许哪些光束进入眼睛。

\subsection{开普勒与伽利略望远镜}

把两薄透镜放在空气中，物镜焦距 \(f_o>0\)，目镜焦距带符号记为 \(f_e\)。远物的像在物镜后焦平面；望远镜调至出射平行光时，该平面同时是目镜的物方焦平面。
\[
 \boxed{d=f_o+f_e},\qquad \boxed{M=-\frac{f_o}{f_e}}.
\]
开普勒式 \(f_e>0\)，两镜间形成真实中间像，\(M<0\)。伽利略式 \(f_e<0\)，目镜在物镜将要会聚的位置之前接收会聚光，\(M>0\)，通常要求 \(f_o>|f_e|\)。这里 \(d\) 是两镜间距。
\fig{o4:kepler}{开普勒式的真实中间像在两镜之间；图以 $f_o:f_e=3:1$ 示意两束平行光的倾角反向}{\begin{tikzpicture}

\draw[axisline](-1.7,0)--(6.2,0);\draw[<->,>=Stealth,world](0,-.8)--(0,1.2);\draw[<->,>=Stealth,world](4,-.8)--(4,1.2);\draw[normal](3,-.9)--(3,1.1);\draw[ray,blue](-1.3,.77)--(0,.9);\draw[ray,blue](0,.9)--(4,.1);\draw[ray,red](4,.1)--(6,-.5);\draw[ray,blue](-1.3,-.43)--(0,-.3);\draw[ray,blue](0,-.3)--(4,.5);\draw[ray,red](4,.5)--(6,-.1);\fill(3,.3)circle(1.5pt)node[above left]{$B_1$};\node[above]at(0,1.2){物镜};\node[above]at(4,1.2){正目镜};\node at(-1,1.5){$\alpha>0$};\node at(5.3,.9){$\alpha'<0$};\draw[dim](0,-1.4)--(4,-1.4)node[midway,below]{$d=f_o+f_e$};

\end{tikzpicture}}
图\ref{o4:kepler}中，同一远物点的两条入射光都到达 $B_1$，随后从目镜前焦平面发散到目镜，并被折成同倾角的平行束。图\ref{fig:o4galileo}则在到达同一会聚目标之前放入负目镜，两者都无焦，但中间光束是否已经交叉不同。焦距比的负号应由这两幅实际光路确定。


\fig{fig:o4galileo}{伽利略望远镜：虚线给出撤去目镜时的会聚目标 \(B_1\)}{
\begin{tikzpicture}[x=1.1cm,y=1.1cm]
\draw[axisline] (-1.5,0)--(4.4,0);
\draw[<->,>=Stealth,world] (0,-.8)--(0,1.25);
\draw[>-<,>=Stealth,world] (2,-.6)--(2,1.3);
\draw[ray] (-1.3,.77)--(0,.9);\draw[ray] (0,.9)--(2,.5);
\draw[ray] (-1.3,-.43)--(0,-.3);\draw[ray] (0,-.3)--(2,.1);
\draw[virtual] (2,.5)--(3,.3);\draw[virtual] (2,.1)--(3,.3);
\draw[ray] (2,.5)--(4.1,1.13);\draw[ray] (2,.1)--(4.1,.73);
\fill (3,.3) circle(1pt) node[below right] {$B_1$};
\node[above] at(0,1.35){物镜};\node[above] at(2,1.4){负目镜};
\draw[dim] (0,-1.15)--(2,-1.15) node[midway,below] {$d=f_o-|f_e|$};
\end{tikzpicture}}
图\ref{fig:o4galileo}中，入射束和出射束分别平行；出射倾角与入射倾角同号。开普勒式的光路和放大率推导见例\ref{ex:o49}。倒像的开普勒望远镜也可加转像系统用于地面观察，因此“倒像”不等于它只能用于天文观察。

望远镜的“无焦”只规定末束仍为平行光，光束方向和直径仍可改变。两类望远镜对中间像的使用也不同：开普勒式让光先完成会聚，伽利略式则在会聚完成前用负目镜接收，所以后者的目镜面对虚物。

当远物转过小角 $\alpha$，物镜中间像高约为 $f_o\alpha$；目镜将该高度变为大小约 $|f_o\alpha/f_e|$ 的出射倾角。再按实际正倒关系赋号，得到 $M=-f_o/f_e$。镜间距条件保证共同焦平面恰好衔接；若间距改变，末束通常不再平行，眼睛就需要重新调节。

\subsection{孔径光阑与出瞳}
望远镜的视角放大率描述方向变化；出瞳描述能进入眼睛的光束截面。出瞳是孔径光阑经其后光学元件所成的像。眼瞳应靠近出瞳，才能接收设计视场内的光束；仅满足焦距比并不能保证观察亮度和视场。

\fig{fv:o4:pupil}{目镜对孔径光阑成像，确定眼睛应放置的轴向位置}{\begin{tikzpicture}[x=.8cm,y=.85cm]
\draw[->](-.5,0)--(9,0);\draw[thick](0,-1.3)--(0,1.3);\draw[thick](5,-1.6)--(5,1.6);\draw[thick](8,-.6)--(8,.6);
\draw[->,blue](0,1)--(5,1)--(8,-.6);\draw[->,red](0,1)--(5,0)--(8,-.6);
\node[above]at(0,1.4){物镜口径};\node[above]at(5,1.7){目镜};\node[above]at(8,.9){出瞳};
\draw[<->](0,-1.9)--(5,-1.9)node[midway,below]{$d$};\draw[<->](5,-1.9)--(8,-1.9)node[midway,below]{$v_p$};
\end{tikzpicture}}
以开普勒望远镜为例，若物镜边缘确为限制光束的光阑，图\ref{fv:o4:pupil}中它距目镜 $d$，出瞳便是目镜对这个圆孔的实像。由 $1/d+1/v_p=1/f_e$，有
\[
 v_p=\frac{df_e}{d-f_e},\qquad D_p=\frac{v_p}{d}D_o.
\]
这里 $D_o,D_p$分别为物镜孔径和出瞳直径；把无穷远调焦条件 $d=f_o+f_e$代入，得到 $D_p=D_o f_e/f_o$。眼睛若偏离这一截面，视场边缘的光束会先被截去。若另有镜筒口径成为限制光阑，就应先找真正的孔径光阑，再求其像，不能继续按物镜口径直接计算。


对伽利略式，若仍由物镜限制孔径，$f_e<0$ 使 $v_p=df_e/(d-f_e)<0$，出瞳为位于目镜前的虚像。镜外眼睛不能实际放在这个虚出瞳平面上，因此必须沿实际出射束检查通光，不能照搬开普勒式“眼睛放在实出瞳”的空间布置。

图\ref{o4:pupil-shift}进一步说明出瞳位置与视场的联系。对无焦系统，同一远物点的出射束平行；若眼瞳离出瞳平面轴向距离为 $z$，出射倾角为 $\alpha'$ 的束中心偏移约 $z\alpha'$。设出瞳半径为 $R_p$、眼瞳半径为 $a$，要完整接收该束，须有
\[
|z\alpha'|+R_p\le a.
\]
\fig{o4:pupil-shift}{出瞳后自由传播的轴上束与斜束；眼瞳偏离出瞳平面后，斜束的部分边缘可被截去}{\begin{tikzpicture}

\draw[axisline](-.7,0)--(4.8,0);\draw[normal](0,-1.4)--(0,1.8);\draw[world](0,-.6)--(0,.6);\fill[gray!30](2.95,.8)rectangle(3.12,1.9);\fill[gray!30](2.95,-1.5)rectangle(3.12,-.8);\draw[world](3,-1.5)--(3,-.8) (3,.8)--(3,1.9);
\foreach\y in {-.6,.6}{\draw[ray,blue](0,\y)--(4.5,\y);}\draw[ray,red](0,-.6)--(4.5,.75);\draw[ray,red](0,.6)--(3,1.5);\node[above]at(0,1.8){出瞳平面};\node[right]at(3.2,1.6){斜束被挡};\node[above]at(3,2.1){眼瞳};\draw[dim](-.5,-.6)--(-.5,.6)node[midway,left]{$2R_p$};\draw[dim](0,-2)--(3,-2)node[midway,below]{$z$};\node at(1.5,1.8){$\alpha'>0$};\draw[dim](5.1,-.8)--(5.1,.8)node[midway,right]{$2a$};

\end{tikzpicture}}
这个不等式描述完整通光的几何条件；不满足时仍可能接收部分光。若 $a<R_p$，轴上束本身也会被眼瞳截断。像的位置、视角放大率、眼睛调焦和孔径截光因而是四项需要分别计算、再合并到同一观察条件中的量。


\subsection{显微镜的两级放大}

物镜先把 \(AB\) 成为放大的实像 \(A_1B_1\)，目镜再把它作为放大镜的物。设物镜焦距 \(f_o\)，目镜焦距 \(f_e\)，两镜间距 \(d\)，光学间隔为 \(\Delta=d-f_o-f_e\)。常规显微镜取 \(\Delta\gg f_o\)。眼睛放松、末像在无穷远时，\(u_e=f_e\)、\(v_o=d-f_e\)，于是
\[
 \beta_o=-\frac{v_o}{u_o}=-\frac{v_o-f_o}{f_o}=-\frac{\Delta}{f_o},
 \qquad \boxed{M_\infty=-\frac{\Delta D}{f_of_e}}.
\]
\fig{fig:o4microscope}{物镜形成倒立实像，目镜把来自同一像点的光折成平行束；图示为放松观察}{
\begin{tikzpicture}[x=1.1cm,y=1cm]
\draw[axisline] (-1.6,0)--(6,0);
\draw[<->,>=Stealth,world] (0,-1.2)--(0,.8);
\draw[<->,>=Stealth,world] (4,-1.2)--(4,.8);
\draw[axis,thick] (-1.090909,0)--(-1.090909,.218182) node[above] {$B$};
\node[below]at(-1.090909,0){$A$};
\draw[axis,thick] (3,0)--(3,-.6) node[below] {$B_1$};\node[above]at(3,0){$A_1=F_e$};
\draw[ray] (-1.090909,.218182)--(0,.218182);
\draw[ray] (0,.218182)--(4,-.872727);\draw[ray] (4,-.872727)--(5.7,.147273);
\draw[ray] (-1.090909,.218182)--(0,0);
\draw[ray] (0,0)--(4,-.8);\draw[ray] (4,-.8)--(5.7,.22);
\node[above]at(0,.9){物镜};\node[above]at(4,.9){目镜};
\node[below]at(.8,0){$F'_o$};
\draw[dim] (.8,-1.5)--(3,-1.5) node[midway,below] {$\Delta$};
\end{tikzpicture}}
图\ref{fig:o4microscope}中 \(A_1B_1\) 位于目镜前焦平面，故出射为平行光。上述 $\beta_o=-\Delta/f_o$ 在薄透镜近轴模型内由 $u_o=f_ov_o/(v_o-f_o)$ 直接得到，不要求再以镜筒长度代替光学间隔。若眼睛贴近目镜、实际近点等于参照距离 $D$，并把末像改在近点，则
\[
 u_e=\frac{f_eD}{D+f_e},\quad v_o=d-u_e,\quad
 M_D=\beta_o\left(1+\frac D{f_e}\right),\quad
 \beta_o=-\frac{v_o-f_o}{f_o}.
\]
不能在这一条件下仍把 \(u_e\) 精确取为 \(f_e\)。常用的 \(|M|\approx dD/(f_of_e)\) 还要求两焦距远小于镜筒长度，因而可用 \(d\) 近似光学间隔。

图\ref{fig:o4microscope}中的两级放大承担不同任务。物镜的横向放大把细小物体变成可作为目镜物的中间像，目镜的视角放大再把这一像转成眼可观察的角度。把物镜的线放大率与目镜的线放大率直接相乘，只能得到末像的长度关系，不能自动得到对眼睛的视角放大率。

若要求改变末像位置后保持同一物体清晰，应先由观察条件确定目镜物距，再反推中间像距和物镜物距。光学间隔 $\Delta$ 是物镜像方焦点到目镜物方焦点的距离，镜筒长度 $d$ 则是两镜间距，两者相差两个焦距。短镜筒或大焦距时，这一差值不能略去；近似公式应在精确衔接关系之后使用。

\subsection{固定镜距下的观察条件与调焦}

若眼睛与目镜有间隔 $\ell$，实际近点为 $D_N$，则最终虚像到目镜的距离大小为 $q=D_N-\ell>0$。应先由目镜确定
\[
u_e=\frac{f_eq}{q+f_e},\qquad v_o=d-u_e,\qquad
u_o=\frac{f_ov_o}{v_o-f_o},
\]
再求物镜横向放大率和总视角放大率
\[
\beta_o=-\frac{v_o-f_o}{f_o},\qquad
M_N=\beta_o\frac{D}{D_N}\left(1+\frac q{f_e}\right).
\]
这样，眼位、近点和固定镜距的约束都落在实际物像位置上。即使只改变最终像的观察距离，也会改变目镜物距，再改变物镜像距和物距；不能在严格计算中只替换目镜放大率，却保留原来的物镜放大率。例\ref{ex:r42}用具体数据比较这两次调焦。

\section{几何光学的尺度与强度边界}

几何光学用光线描述传播方向，在波长远小于相关孔径和几何尺度时能够近似计算成像位置。孔径的衍射会限制分辨能力；即使几何光线都指向同一点，实际焦斑也不会无限小。光线图的会聚位置与实际能分辨的细节，应分别判断。

若透光孔径的特征尺度为 $a$，波长为 $\lambda$，几何光学以 $\lambda/a\ll1$ 为小参数。图\ref{o4:finite-spot}区分理想光线交点与有限口径下的焦斑宽度。忽略像差时，衍射给出的角宽尺度约为 $\lambda/a$，焦平面宽度尺度约为
\[
\Delta x\sim\frac{f\lambda}{a}.
\]
比例系数与孔径形状、所用宽度定义有关，这里只作尺度估计。
\fig{o4:finite-spot}{光线交点确定几何像的位置，有限孔径的波动传播决定焦斑宽度；右图只示意宽度，不表示具体强度函数}{\begin{tikzpicture}

\begin{scope}[scale=.85]
\draw[axisline](-1.8,0)--(3.6,0);\draw[<->,>=Stealth,world](0,-1)--(0,1);\foreach\y in {-.7,.7}{\draw[ray,blue](-1.6,\y)--(0,\y);\draw[ray,blue](0,\y)--(3,0);}\fill(3,0)circle(1.5pt)node[below]{$F$};\node at(.8,1.7){几何像点};\end{scope}
\begin{scope}[xshift=5.3cm,scale=.85]
\draw[axisline](-1.8,0)--(3.6,0);\draw[<->,>=Stealth,world](0,-1)--(0,1);\foreach\y in {-.7,.7}{\draw[ray,blue](-1.6,\y)--(0,\y);\draw[guide](0,\y)--(3,0);}\fill[red!25](3,0)ellipse(.12 and .22);\draw[world,red](3,0)ellipse(.12 and .22);\draw[dim](3.5,-.22)--(3.5,.22);\node[right]at(3.5,.5){$\Delta x$};\draw[dim](-.5,-.7)--(-.5,.7)node[midway,left]{$a$};\draw[dim](0,-1.3)--(3,-1.3)node[midway,below]{$f$};\node at(.8,1.7){有限孔径焦斑};\end{scope}

\end{tikzpicture}}
当所问的是像的大致位置，这个宽度可以相对系统尺寸很小；当所问的是邻近细节能否分开，它却可能决定结果。放大镜或目镜只扩大观察视角，不能把已经混在同一焦斑中的信息重新分离。第五讲将把具体路径的相位差和强度叠加写入波动计算。


用经典电磁场描述时，干涉和衍射来自场振幅的相干叠加；用单光子描述时，则计算相应的概率幅。单个光子的探测结果是局域事件，反复积累仍可形成干涉条纹；因此不能简单按光子数多少把波动性和粒子性分开。光电效应中，单次吸收能量 \(h\nu\) 进入电子与物质系统的能量平衡。光既不等同于经典小球，也不要求用机械介质承载其在真空中的传播。

用反射镜和透镜产生空心光束时，光路设计采用几何光学的细束近似。若要计算中心暗区的实际光强、衍射填充和传播距离，就须进一步给出波长、入射场分布、各元件孔径，并用波动光学计算。几何设计保证光线的传播方向，不单独决定全部强度分布。
